% answers_short.tex —— 2007 年·北京卷·文科 · 简短答案速查
% 由 tools/make_short_answers.py 自 answers.tex 的 \gksolhead/\gkanswer 提取，
% 并与 2007/2007_数学真题_北京文_答案要点.md 逐题比对后生成，勿手改（改源 MD 后重跑）。

\begin{gknotice}
\textbf{使用说明}：本节只给最终答案便于快速核对，完整推导见解析版。答案为编写者独立推导，\textbf{非官方标准答案}。
\end{gknotice}

\gkpart{选择题}
\begin{center}\begin{tabular}{ccccccccc}
\toprule
题号 & 1 & 2 & 3 & 4 & 5 & 6 & 7 & 8 \\
\midrule
答案 & C & B & B & D & A & C & D & C \\
\bottomrule
\end{tabular}\end{center}

\gkpart{填空题}
\begin{itemize}[leftmargin=2.6em]
  \item[\textbf{9.}] $3$
  \item[\textbf{10.}] $a_n=2n-11$
  \item[\textbf{11.}] $\lambda=-3$
  \item[\textbf{12.}] $\dfrac{\sqrt{10}}2$
  \item[\textbf{13.}] $\dfrac7{25}$
  \item[\textbf{14.}] $1$；$x=1$
  \item[\textbf{15.}] (1) $P=(-1,3)$；(2) $a>2$
  \item[\textbf{16.}] (1) $c=2$；(2) $a_n=n^2-n+2$
\end{itemize}

\gkpart{解答题}
\begin{itemize}[leftmargin=2.6em]
  \item[\textbf{17.}] (1) 见解析版；(2) $\arccos\dfrac{\sqrt6}{4}$
  \item[\textbf{18.}] (1) U/PENDING；若进一步假设六位乘客的选择相互独立，则概率为 $\dfrac{10\cdot9\cdot8\cdot7\cdot6\cdot5}{10^6}=\dfrac{189}{1250}$；(2) U/PENDING；若进一步假设六位乘客的选择相互独立，则概率为 $\binom63(1/10)^3(9/10)^3=\dfrac{729}{50000}$
  \item[\textbf{19.}] (1) $3x+y+2=0$；(2) $(x-2)^2+y^2=8$；(3) $\dfrac{x^2}{2}-\dfrac{y^2}{2}=1$ 的左支（$x\leq-\sqrt2$）
  \item[\textbf{20.}] (1) $k>2\sqrt2$；(2) $t=\dfrac{x_1}{2}-\dfrac1{x_1}$；定义域 $(0,\sqrt2)$，值域 $(-\infty,0)$；(3) $|OM|=|ON|$
\end{itemize}
